% Generated draft: requires mathematical and source review.
\documentclass[11pt]{article}
\usepackage[margin=25mm]{geometry}
\usepackage{fontspec}
\setmainfont{Latin Modern Roman}
\usepackage{amsmath,amssymb,mathtools,unicode-math}
\setmathfont{Latin Modern Math}
\usepackage{enumitem,needspace,xurl,hyperref}
\setlength{\parindent}{0pt}
\setlength{\parskip}{5pt}
\setlength{\emergencystretch}{3em}
\newcommand{\N}{\mathbb{N}}
\newcommand{\Z}{\mathbb{Z}}
\newcommand{\Q}{\mathbb{Q}}
\newcommand{\R}{\mathbb{R}}
\newcommand{\C}{\mathbb{C}}
\newcommand{\F}{\mathbb{F}}
\newcommand{\A}{\mathbb{A}}
\newcommand{\PP}{\mathbb{P}}
\newcommand{\E}{\mathbb{E}}
\newcommand{\eps}{\varepsilon}
\newcommand{\dd}{\,\mathrm{d}}
\begin{document}
\section*{UnsolvedMath ID: 30003876}
\textbf{Source problem code:} catalogue entry
\subsection{Definitions and mathematical statement}
Let \(n\geq 6\) be an integer, and let \(A_n\) denote the alternating group on \(n\) letters. Let \(g\) and an embedding
\[
\iota\colon A_n\hookrightarrow \operatorname{Sp}_{2g}(\F_2)
\]
be as stipulated in the supplied record. Here \(\operatorname{Sp}_{2g}(\F_2)\) is the group preserving a nondegenerate alternating form on \(\F_2^{2g}\). Choose a Siegel parabolic, meaning the stabilizer of a Lagrangian subspace, and let \(U\) denote the group of \(\F_2\)-points of its unipotent radical. Identify \(A_n\) with its image under \(\iota\), and assume that
\[
E=A_n\cap U
\]
is a maximal elementary abelian \(2\)-subgroup of \(A_n\). Thus \(E\cong (\Z/2\Z)^e\) for some integer \(e\geq 0\), and it is not properly contained in another elementary abelian \(2\)-subgroup of \(A_n\). Define
\[
e_{A_n}=\dim_{\F_2}E.
\]
Maximality here means maximality under inclusion; no assertion that \(E\) has the largest possible rank is added.

In the intended moduli interpretation, \(A_g\) denotes the moduli space of principally polarized abelian varieties of dimension \(g\). Use the cover specified in the record,
\[
\widetilde A_{g,2}\longrightarrow A_g,
\]
described there as the Galois closure of the universal \(2\)-torsion cover, with the intended symplectic group action. Define the intermediate quotient
\[
A_g^{A_n}:=\widetilde A_{g,2}/\iota(A_n).
\]
The record does not specify the ground field, whether moduli spaces or stacks are intended, or the precise universal \(2\)-torsion cover and its Galois action. These choices must be fixed before this notation determines a particular cover.

For clarity, under the usual convention for integral varieties over a specified field \(k\), the resolvent degree \(\operatorname{RD}_k(X\to Y)\) of a dominant generically finite rational map is the least nonnegative integer \(d\) for which there exists a tower of dominant generically finite rational maps
\[
Y_r\longrightarrow Y_{r-1}\longrightarrow\cdots\longrightarrow Y_0=Y
\]
whose terminal variety admits a dominant rational map to \(X\) over \(Y\), and in which each step is birational to an irreducible component of a rational pullback of a dominant generically finite map between varieties of dimension at most \(d\). The supplied notation \(\operatorname{RD}\) leaves the ground field and any adaptation to moduli stacks implicit.

\textbf{Conjectural inequality.} For the data and intended cover above, is
\[
\operatorname{RD}\!\left(\widetilde A_{g,2}\longrightarrow A_g^{A_n}\right)\geq e_{A_n}?
\]
The original statement restricts the embedding to one described in the report, whereas the cleaned statement gives the intersection condition above. The supplied record does not establish whether that condition exhausts the restrictions on the intended embeddings or genera. This draft therefore records the inequality subject to those unresolved qualifications, without asserting existence of the required embeddings for specified values of \(g\).
\subsection{Short English statement}
For \(n\geq 6\), embed the alternating group \(A_n\) into a finite symplectic group as prescribed by the source. Suppose its intersection with the chosen Siegel unipotent radical is a maximal elementary abelian \(2\)-subgroup of rank \(e_{A_n}\). The question asks whether the corresponding quotient cover of the moduli space has resolvent degree at least \(e_{A_n}\). The ground field, precise moduli construction, and full restrictions on the embedding require clarification. The supplied record labels the problem open; this draft does not independently verify that status.
\subsection{Sources}
\begin{itemize}
\item Local UnsolvedMath registry, permanent ID \(30003876\), title \emph{Lower Bounds for Symplectic Resolvent Degrees}, source problem code \texttt{catalogue entry}, in the supplied Oberwolfach Reports Open Problems collection. No individual proposer is identified; the proposed year is recorded as 2018.
\item Supplied source citation: \emph{Algebraische Zahlentheorie (2019)}, Oberwolfach Reports, pp.\ 1785--1844, DOI \texttt{10.4171/owr/2018/30}. \url{https://doi.org/10.4171/owr/2018/30}. The citation's displayed year is retained as supplied.
\item Supplied literature source: \emph{Oberwolfach Report 30/2018}. \url{https://ems.press/content/serial-article-files/46751}.
\item Supplied literature source: Farb--Wolfson, \emph{resolvent degree}. \url{https://arxiv.org/abs/1803.04063}. The record supplies no more detailed bibliographic citation for this reference.
\item Literature assessment dated exactly \texttt{2026-08-21}: supplied dated curation, not independently verified. It labels the problem open and describes an unproved lower bound involving the exponent of an abelianization, together with an unestablished implied Hilbert conjecture. That description does not directly match the rank-based inequality displayed in this record, and its relationship to the present target requires review.
\item The registry supplies the statement status \texttt{corrected\_verified} and reports a check against the official Oberwolfach report. These are supplied provenance claims; this draft does not independently verify the report, the literature assessment, or the mathematical formulation.
\end{itemize}
\end{document}
